\section{Reverse Reachability Has No Domain Here}
\label{sec:inapplicability}

We give two independent obstructions. The first is a base-case failure and is easy to verify
empirically; the second is more fundamental and rules out any repair of the reduction.

\subsection{The reverse walk has no base case}

\begin{proposition}[Null activation threshold]
\label{prop:null}
Let $(\theta^\kappa_v, \theta^\tau_v)$ be drawn with density $2$ on the simplex
$\Delta$ of Eq.~\eqref{eq:simplex}. Then $\Pr[\theta^\kappa_v = 0] = 0$, and consequently a node
with $\delta(v,t) = 0$ is positively activated with probability $g(0) = 0$.
\end{proposition}

\begin{proof}
$\{\theta^\kappa = 0\}$ is the lower-dimensional face of $\Delta$ and has measure zero under a
density on the two-dimensional simplex. The second statement is Remark~\ref{rem:lemma1}:
$g(0) = 2 \cdot 0 \cdot 1 = 0$.
\end{proof}

Now consider the reverse walk used to build $\mathcal{R}$. Such a walk traverses edges backwards
from a source and admits a predecessor $u$ into $\mathcal{R}$ when $u$ \emph{could be active given
that it has received no influence}. Under Eq.~\eqref{eq:transition} this requires
$\theta^\kappa_u \le 0 \le \theta^\tau_u$, hence $\theta^\kappa_u = 0$. By
Proposition~\ref{prop:null} that event has probability zero, so no predecessor is ever admitted
and the walk terminates immediately.

\begin{corollary}[Empty RR sets]
\label{cor:empty}
For a non-seed source, $\mathcal{R} = \emptyset$ almost surely. The identity
Eq.~\eqref{eq:rr-identity} therefore returns $0$ for every seed set, against a true spread that is
generally positive.
\end{corollary}

\paragraph{Empirical confirmation.}
Table~\ref{tab:rr} reports the measurement on NetHEPT ($n = 15233$, $m = 32235$, in-weights
normalised so that $\sum_u \omega_{uv} = 1$ for every $v$ with in-edges, matching the source
model's setting). We ran $300$ reverse breadth-first walks from uniformly random sources, where a
predecessor is admitted exactly when $0 \in [\theta^\kappa_u, \theta^\tau_u]$.

\begin{table}[t]
\centering
\caption{Reverse reachability on NetHEPT. Every RR set is empty, so the classical estimator
degenerates to zero.}
\label{tab:rr}
\begin{tabular}{lr}
\toprule
quantity & value \\
\midrule
nodes whose window contains $0$ & $0 \;/\; 15233$ \\
RR-set size (min / median / mean / max) & $0 \;/\; 0 \;/\; \mathbf{0.0} \;/\; 0$ \\
RR-set size as fraction of the graph & $0.0000$ \\
TIM-style estimate $n \cdot \Pr[S \cap \mathcal{R} \neq \emptyset]$, $|S| = 50$ & $\mathbf{0.0}$ \\
true Monte-Carlo spread $\sigma(S)$, $|S| = 50$ & $111.93$ \\
\bottomrule
\end{tabular}
\end{table}

\subsection{The reduction has no randomised structure to sample}

The base-case failure above is about the simplex draw. One might hope to repair it, for example
by placing mass at $\theta^\kappa = 0$, or by treating seeds as the only permitted base case.
The second obstruction shows why such repairs do not restore Eq.~\eqref{eq:rr-identity}.

\begin{proposition}[No edge-factorised structure]
\label{prop:nofactor}
Under the process of \eqref{eq:delta}--\eqref{eq:transition}, the event ``$v$ is ever positively
activated'' is not a function of a per-edge independent random structure. Consequently no
live-edge graph and no trigger set exists whose reachability relation reproduces the process.
\end{proposition}

\begin{proof}[Proof sketch]
By Observation~\ref{obs:setfunction}, after conditioning on the window draw the entire trajectory
--- including which nodes are ever positive and the final spread --- is a deterministic function
of $S$. The only random object is the window draw, which is indexed by \emph{node}, not by edge,
and which does not fail independently per edge. The activation rule \eqref{eq:transition} depends
on the \emph{sum} $\delta(v,t)$ of weights over a dynamically growing set
$A^{\mathrm{in}}_t(v)$; it is therefore not expressible as ``at least one incident edge
succeeded'' (the IC form) nor as ``the set of active in-neighbours covers a sampled threshold
set'' with a per-edge independent threshold (the LT form). Since the reachability relation in
Eq.~\eqref{eq:rr-identity} must be induced by such a structure, and no such structure exists, the
identity is unavailable.
\end{proof}

Two consequences deserve emphasis.

\begin{itemize}
  \item \textbf{This is a domain problem, not an efficiency problem.} The RR family --- TIM,
        TIM+, IMM, OPIM, SSA, D-SSA --- and its GPU-parallel descendants all consume the same
        identity. None of them can be instantiated, so their asymptotic guarantees and their
        orders-of-magnitude speedups do not transfer. There is no RR baseline to beat in this
        setting.
  \item \textbf{The result does not depend on non-monotonicity.} Even if $g$ were replaced by a
        monotone function, Proposition~\ref{prop:nofactor} would still hold, because it concerns
        the factorisation of the randomness rather than the shape of the objective. This matters
        because the objective in this setting turns out to be monotone over most of the tested
        range (\S\ref{sec:regime}), so an argument phrased purely in terms of non-monotonicity
        would be fragile.
\end{itemize}

\subsection{The source model avoids the objective, and we can say how}

It is worth recording that the model's own solution path is consistent with the above. Its
analysis establishes only a $\gamma/k$ ratio for incremental greedy --- not $(1-1/e)$ --- and its
upper-bound method is constructed by exhibiting two \emph{monotone submodular} functions
$\sigma^\kappa(\cdot)$ and $\sigma^\tau(\cdot)$ derived from linear-threshold processes, with
$\lambda(S) \ge \sigma(S)$, and then optimising the surrogate $\lambda$. Since
$\sigma^\kappa$ and $\sigma^\tau$ are genuine LT influence functions, RR sampling \emph{can} be
applied to them. The correct statement is therefore not ``RR is impossible in this problem'' but:

\begin{quote}
\emph{RR can only be applied to a monotone upper bound $\lambda(S) \ge \sigma(S)$, whose
optimisation guarantees a $\gamma/k$ ratio at best; applying coverage semantics directly to
$\sigma$ makes the estimator degenerate to zero.}
\end{quote}

Quantifying the surrogate gap $\lambda(S) - \sigma(S)$ is a natural next step and we flag it in
\S\ref{sec:limitations}.
